\chapter{Fejlesztői dokumentáció}
\label{ch:impl}

\section{Program architektúra}

\section{Növényzet generálási paraméterek}


\subsection{Tájfutótérképek}

A tájfutó térképek elsősorban a vegetációt alkotó növényegyedek elszórásához szükséges adatokat tartalmaznak.
Mint már előzőleg említve lett, két fajta térkép szimbólum típust különböztetünk meg: \todo{tényleg megkéne korábban említeni} egy jelölés akadályozhatja vagy teret adhat egy növényzet terjedésének. Ennek eldöntésének érdekében szükségünk van a térképeken feltüntetett objektumok és velük asszociált szimbólumok kinyerésére és feldolgozására. 
\newline

\todo{Kód arról, hogy hogyan szerezzük meg az objektumokat és szimbólumokat}

A programcsomag két különböző fájlformátumot tud értelmezni és feldolgozni.
Az első ilyen formátum az ún. "OCAD" fájlformátum, amely a ".ocd" fájlkiterjesztéssel van leggyakrabban ellátva.
Ez egy elsősorban az OCAD programcsomag által használt és mellett fejlesztett bináris fájl formátum, amely a orientációs térképeken megtalálható szimbólumat, objektumokat, színeket, stb. kódolja el.
Előny ennél a formátumnál az, hogy több nagyobb szoftver csomag is ráépül illetve támogatja, így ennél a szoftvernél való támogatottsága is természetesen adódik.
Nehézsége azonban, hogy az így megformázott fájlok inkább fekete doboz alapon működnek. \todo{Észlelések leírása}

A második támogatott fájlformátum az elsősorban a nyílt-forráskódú "OpenOrienteering Mapper" program által használt ".omap" kiterjesztésü OpenOrienteering fájlformátum. Ez egy XML alapú struktúrában tárolja el a különböző orientációs objektumokat és hozzájuk tartozó szimbólumokat.
Az OCAD formátummal szemben ez egy sokkal átláthatóbb adathordozónak bizonyult a téma kidolgozása alatt, így implementálása és tesztelése is egyszerűbbnek bizonyult. 

\section{Növények szétszórása}

\section{Növények ütközésvizsgálata objektumokkal}

\section{Rosszul formált térképek kezelése}

Erdő mint szimbólum kihagyása

\section{Egyedek fajhoz való besorolása}

\subsection{Kezdeti egyed halmaz megállapítása}

\subsection{Szomszédossági besorolás}

\section{Vegetáció területének maximalizálása}

\section{Növényzet megjelenítés}

Folyamat leírása:

AlgLib könyvtár használata + Nanoflann (KD fa felépítése szomszédok megtalálására)

Részgráfokra bontás -> Mely növények vannak közel egymáshoz, ezek tudják egymást "méretét" korlátozni

Részgráfok felépítése, majd egymás közötti relációk felépítése AlgLib számára

Optimalizáció elvégzése

0 nagyságú növények eltávolítása

Folyamat megismétlése hogy a méreteket stabilizáljuk

Megfontonlandó dolgok:

Az eredmények pontossága -> Meg lehet figyelni, hogy egyes lineáris megkötések, amelyek a növények közötti távolságon alapulnak, nem fognak teljesülni, olyan 0.00001-es nagyságrendűleg pontatlanság figyelhető meg

Okok: Optimalizációs folyamat megállási paraméterein kell fejleszteni -> Milyen paraméterek lennének esetünkben megfelelőek
\todo{}

